\section{Finite-difference penalized trend estimation}

\subsection{Penalized least squares}

At forecast origin $T$, let
\begin{equation}
\label{eq:observation_window}
x_T=(y_{T-L+1},\ldots,y_T)^\top\in\mathbb{R}^{L}
\end{equation}
denote the estimation window. For difference order $d<L$, let
\begin{equation}
\label{eq:penalty_gram_matrix}
D_d\in\mathbb{R}^{(L-d)\times L},
\qquad
Q=D_d^\top D_d.
\end{equation}
The baseline estimator solves
\begin{equation}
\label{eq:pls_problem}
\widehat\tau_T(\lambda)
=
\arg\min_{\tau\in\mathbb{R}^{L}}
\left\{
\|x_T-\tau\|_2^2+\lambda\|D_d\tau\|_2^2
\right\},
\qquad \lambda\ge0,
\end{equation}
with solution
\begin{equation}
\label{eq:H}
\widehat\tau_T(\lambda)
=
H_\lambda x_T,
\qquad
H_\lambda=(I+\lambda Q)^{-1}.
\end{equation}

The paper uses the zero-drift, identity-weight formulation in \cref{eq:pls_problem} to isolate the smoothness-selection question. Guerrero's broader statistical formulation allows nonzero mean differences and more general covariance weighting \cite{guerrero2007penalized}; these are natural extensions but are not silently imposed in the baseline criterion.

\subsection{Spectral representation}

Because
\begin{equation}
\label{eq:penalty_positive_semidefinite}
v^\top Qv=\|D_dv\|_2^2\ge0,
\end{equation}
$Q$ is symmetric positive semidefinite. The dimension of its null
space, and the existence of the smoother at every finite penalty, follow
directly from the difference operator.

\begin{proposition}[Null space and finite-penalty well-posedness]
\label{prop:pls_wellposed}
For the standard consecutive order-$d$ difference operator with
$1\le d<L$,
\begin{equation}
\label{eq:penalty_rank}
\operatorname{rank}(D_d)=\operatorname{rank}(Q)=L-d,
\qquad
\ker(Q)=\ker(D_d),\qquad
\dim\ker(Q)=d.
\end{equation}
For every finite $\lambda\ge0$, $A_\lambda=I+\lambda Q$
and $H_\lambda=A_\lambda^{-1}$ are symmetric positive definite and
invertible. The objective in \cref{eq:pls_problem} therefore has the
unique minimizer $H_\lambda x_T$.
\end{proposition}

\begin{proof}
The condition $D_dv=0$ is the order-$d$ homogeneous
difference recurrence. Its first $d$ entries are arbitrary, and the
remaining entries are uniquely determined because the coefficient of
each successive new entry is nonzero. Thus
$\dim\ker(D_d)=d$, so rank--nullity gives
$\operatorname{rank}(D_d)=L-d$.
Clearly $D_dv=0$ implies $Qv=D_d^\top D_dv=0$.
Conversely, if $Qv=0$, then
$0=v^\top Qv=\|D_dv\|_2^2$, which implies $D_dv=0$.
The kernels are therefore equal, and rank--nullity also gives the
stated rank of $Q$. Finally, for every nonzero $v$,
\begin{equation}
\label{eq:finite_penalty_spd}
v^\top A_\lambda v
=\|v\|_2^2+\lambda\|D_dv\|_2^2>0
\qquad(0\le\lambda<\infty).
\end{equation}
Hence $A_\lambda$ is symmetric positive definite and invertible;
its inverse $H_\lambda$ is likewise positive definite. The Hessian
of the penalized objective is $2A_\lambda\succ0$; the objective is
strictly convex and coercive, so its stationary point $H_\lambda x_T$
is its unique minimizer.
\end{proof}

In particular, $Q$ has exactly $d$ zero eigenvalues and
$L-d$ strictly positive eigenvalues. Write
\begin{equation}
\label{eq:penalty_spectral_decomposition}
Q=
U\operatorname{diag}
\left(
\underbrace{0,\ldots,0}_{d},
\delta_1,\ldots,\delta_{L-d}
\right)U^\top,
\qquad
\delta_j>0.
\end{equation}
Then
\begin{equation}
\label{eq:H_spectral}
H_\lambda
=
U\operatorname{diag}
\left(
\underbrace{1,\ldots,1}_{d},
\frac{1}{1+\lambda\delta_1},
\ldots,
\frac{1}{1+\lambda\delta_{L-d}}
\right)U^\top.
\end{equation}

Thus the $d$ null-space directions are reproduced without shrinkage, while each penalizable direction is attenuated by
\begin{equation}
\label{eq:spectral_shrinkage}
\alpha_j(\lambda)=\frac{1}{1+\lambda\delta_j}.
\end{equation}
For $d=2$, the null space consists of discrete linear sequences, so increasing $\lambda$ moves the fitted path toward a linear trend.

For finite $\lambda$, the $d$ zero eigenvalues of $Q$ become
$d$ unit eigenvalues of $H_\lambda$; the remaining eigenvalues
belong to $(0,1]$. Consequently $\ker(H_\lambda)=\{0\}$ for
finite $\lambda$, despite $\dim\ker(Q)=d$. As $\lambda$ grows,
$A_\lambda$ can become ill-conditioned even though it remains
invertible, motivating spectral shrinkage and exact treatment
of the limiting projection rather than numerical inversion
at an arbitrarily large penalty.

\subsection{Normalized smoothness}

Guerrero's trace-based construction motivates measuring smoothing through the effective action of $H_\lambda$. We use
\begin{equation}
\label{eq:S}
S(\lambda)
=
1-
\frac{1}{L-d}
\sum_{j=1}^{L-d}
\frac{1}{1+\lambda\delta_j}.
\end{equation}

For comparison, the original Guerrero-type index is
$S^{\mathrm{G}}_d(\lambda)=1-\operatorname{tr}(H_\lambda)/L$.
Our full-range normalization is simply
\begin{equation}
\label{eq:smoothness_index_rescaling}
S(\lambda)=\frac{L}{L-d}S^{\mathrm{G}}_d(\lambda),
\qquad d\ge1.
\end{equation}
Thus the rescaling preserves the penalty ordering and every fitted trend;
it should not be interpreted as a new smoothing estimator
\cite{guerrero2007penalized,cortes2017smoothness}.

The normalization excludes the $d$ directions that cannot be penalized, so
\begin{equation}
\label{eq:smoothness_endpoints}
S(0)=0,
\qquad
\lim_{\lambda\to\infty}S(\lambda)=1.
\end{equation}

For a fixed linear smoother,
\begin{equation}
\label{eq:edf_definition}
\operatorname{edf}(\lambda)=\operatorname{tr}(H_\lambda).
\end{equation}
Combining \cref{eq:H_spectral} and \cref{eq:S},
\begin{equation}
\label{eq:edf}
\boxed{
\operatorname{edf}(\lambda)
=
L-(L-d)S(\lambda).
}
\end{equation}
Hence $S$ is a normalized complement of effective degrees of freedom.

The practical CV search uses $S$, not a truncated grid of $\lambda$.
The parameter transformation covers both limiting smoothing regimes:
$S=0$ corresponds to $\lambda=0$, whereas $S=1$ corresponds to
$\lambda\to\infty$, the least-squares projection onto the
$d$-dimensional null space of $D_d$. A uniform grid in $S$
therefore has an interpretable common bounded domain, unlike an
arbitrarily bounded grid in $\lambda$. Because the mapping is
strictly increasing, it does not invent new minima or make the
forecast loss convex; its advantage is coordinate choice and
endpoint handling.
